\section{Objet et statut de l'annexe}
\label{annexe-num-objet}

Cette annexe rassemble les protocoles numériques utilisés dans les
Tomes I et II. Elle ne remplace pas les démonstrations des annexes
précédentes. Elle précise :

\begin{itemize}
\item
  quelles quantités ont été calculées ;
\item
  comment les seuils et les durées ont été repérés ;
\item
  quelles expériences distinguent amplitude, phase et mémoire
  d'interface ;
\item
  quels contrôles empêchent d'attribuer un effet au mauvais
  mécanisme ;
\item
  quelles conclusions sont robustes et quelles questions restent
  ouvertes.
\end{itemize}

Dans cette annexe, trois étiquettes de la taxonomie générale sont particulièrement utilisées.

\begin{center}
\begin{tabular}{p{0.18\linewidth}p{0.70\linewidth}}
\hline
\texttt{[PROP-NUM]}
&
Un calcul analytique définit l'objet, puis une résolution numérique
fournit sa valeur : racine, valeur propre, multiplicateur de Floquet,
période ou seuil.
\\
\texttt{[NUM]}
&
Le résultat provient d'une intégration, d'un balayage de paramètres ou
d'une classification de trajectoires.
\\
\texttt{[CTRL]}
&
Une expérience est construite pour éliminer une interprétation
concurrente.
\\
\hline
\end{tabular}
\end{center}

Les simulations ne constituent pas une validation empirique du modèle.
Elles vérifient sa cohérence interne, testent ses mécanismes et
délimitent les affirmations permises.

\section{Fichiers de calcul et traçabilité}
\label{annexe-num-tracabilite}

Les résultats du cycle événementiel, des bandes de phase et de la
mémoire d'interface ont été produits par les scripts :

\begin{verbatim}
fii11_cycle.py
fii12_bands_specificity.py
\end{verbatim}

Les sorties structurées correspondantes sont :

\begin{verbatim}
output_theory/fii11_report.json
output_theory/fii12_report.json
\end{verbatim}

Les vérifications complémentaires, notamment le prédicteur de signe,
le contrôle d'habituation douce et la comparaison des mémoires
scalaire et vectorielle, sont consignées dans l'historique du dépôt.

Le paquet de reproductibilité associé à cette version du manuscrit est
accessible dans le dépôt public
\url{https://github.com/bhurpeau/IWS_2}. La révision de référence est le
commit \href{https://github.com/bhurpeau/IWS_2/commit/a39c2d2b6284964134fdb1052e4556ca6ae5eab3}{\texttt{a39c2d2}}
(\emph{New repo IWSV2}, 1er septembre 2026). Cette empreinte, plutôt que
la branche mobile \texttt{main}, fixe le contenu exact auquel se
rapportent les scripts et sorties mentionnés dans la présente annexe.

Pour qu'un résultat numérique soit reproductible, une archive de
référence doit conserver ensemble :

\begin{itemize}
\item le code exact ;
\item les paramètres du modèle ;
\item les conditions initiales ;
\item le protocole d'exposition ;
\item le maillage des valeurs testées ;
\item les tolérances du solveur et de la détection d'événements ;
\item la règle de classification finale ;
\item le rapport produit par le script.
\end{itemize}

Cette annexe fixe les éléments conceptuels et les valeurs utilisées.
Les réglages de bas niveau du solveur doivent rester ceux déclarés dans
la version archivée des scripts ; ils ne sont pas réinventés ici.

\section{Réduction et conventions numériques}
\label{annexe-num-conventions}

Sauf mention contraire, les expériences du Tome II utilisent la
réduction diagonale

\[
h=(x,x),
\qquad
v=(\dot x,\dot x),
\qquad
\tau=(\theta,\theta).
\]

Le système source complet est maintenu sur son puits positif, de sorte
que la projection transmise par composante est

\[
w=\tanh(x_A)=0{,}98006>0.
\]
Cette valeur est celle mesurée sur l'attracteur source utilisé par le
protocole ; l'écriture \(x_A\simeq2{,}30\) n'en est qu'un arrondi de
présentation.

\subsection*{Paramètres de référence}

Sauf mention locale contraire, le jeu de référence utilisé pour le
noyau et sa réduction bidimensionnelle est le suivant. Les valeurs du
mécanisme de Kairos sont celles déclarées dans la configuration du
modèle canonique utilisée par les programmes de référence.

\begin{center}
\begin{tabular}{p{0.17\linewidth}p{0.18\linewidth}p{0.52\linewidth}}
\hline
Paramètre & Valeur & Rôle \\
\hline
$d$ & $2$ & dimension de la réduction de référence \\
$\zeta$ & $0{,}8$ & gain du couplage relationnel \\
$\lambda$ & $2$ & anisotropie du préconditionneur historique \\
$\alpha$ & $0{,}45$ & taux d'écriture interne de la trace \\
$\beta$ & $0{,}12$ & taux d'oubli de la trace \\
$\gamma_0$ & $0{,}25$ & friction minimale \\
$\gamma_1$ & $1{,}25$ & dépendance de la friction à la trace \\
$\kappa_1$ & $0{,}6$ & contribution de $\lVert H\rVert$ à la pression \\
$\kappa_2$ & $0{,}4$ & contribution de $\lVert V\rVert$ à la pression \\
$\kappa_3$ & $0{,}5$ & relaxation de la pression \\
$\Theta$ & $1{,}25$ & seuil de déclenchement des Kairos \\
$r_P$ & $0{,}35$ & relaxation de la pression au saut \\
$r_V$ & $0{,}5$ & relaxation de la vitesse au saut \\
\hline
\end{tabular}
\end{center}

Les paramètres propres aux interfaces ne font pas partie de ce socle.
Pour les expériences correspondantes, les jeux utilisés sont

\[
\text{R1:}\qquad \rho=0{,}05,\qquad \eta=1,
\]

\[
\text{R2:}\qquad
\lambda_\chi=0{,}2,\qquad
\kappa_Q=0{,}6,\qquad
r_\chi=0{,}8,\qquad
\kappa_\chi=2{,}5.
\]

L'extension à mémoire lente de l'interface utilise, lorsqu'elle est
activée,
\[
\varepsilon_s=0{,}005,
\qquad
c_s=8.
\]
Ces valeurs sont des réglages des protocoles d'interface étudiés et ne
modifient pas le jeu canonique du noyau.

Les deux états autonomes principalement distingués après découplage
sont :

\[
x\simeq0
\]

pour le repos, et

\[
x\simeq+2{,}30
\]

pour le puits positif.

La classification d'une interaction ne repose pas sur l'état au moment
du découplage. Après l'exposition, la trajectoire est poursuivie sous
la dynamique autonome. Il y a appropriation seulement si elle rejoint
un attracteur autonome différent de celui associé à la condition
initiale.

\section{Détection des événements}
\label{annexe-num-detection-evenements}

Les Kairos sont définis par le franchissement

\[
p(t)=\Theta.
\]

Une simulation hybride doit donc :

\begin{enumerate}
\item intégrer le flot jusqu'au premier franchissement ;
\item localiser l'instant de l'événement ;
\item appliquer la loi de saut ;
\item redémarrer l'intégration depuis l'état post-événement ;
\item traiter sans avancement de temps toute cascade simultanée.
\end{enumerate}

Une simple discrétisation par pas fixes sans localisation du seuil
déforme la période \(\Delta\), la charge pré-événement \(\chi^-\), le
quantum injecté dans la trace et la phase à laquelle une exposition
ultérieure rencontre le système.

Cette précision est particulièrement importante pour les bandes de
phase : un faible décalage des événements peut déplacer les bords des
bandes sans modifier leur mécanisme.

\section{Continuation de la branche résistante}
\label{annexe-num-continuation-resistance}

La branche résistante du modèle continu R0 est définie par le résidu

\[
F_g(x)=0.
\]

La continuation numérique suit la racine négative issue de la
bifurcation transcritique en

\[
g_c=0{,}0327.
\]

Pour chaque valeur de \(g\) :

\begin{enumerate}
\item la racine précédente sert d'initialisation à la résolution suivante ;
\item l'équilibre complet \((x_{\mathrm{res}},0,\theta_{\mathrm{res}})\) est reconstruit ;
\item le Jacobien tridimensionnel est évalué ;
\item ses trois valeurs propres sont calculées ;
\item la plus grande partie réelle est enregistrée.
\end{enumerate}

Le balayage de référence contient \(281\) valeurs de \(g\). La paire
complexe dominante vérifie :

\[
\max\Re\operatorname{spec}(J)
=
-0{,}0025
\quad
\text{à }g=0{,}120,
\]

\[
\max\Re\operatorname{spec}(J)
=
0
\quad
\text{à }g=0{,}1285,
\]

\[
\max\Re\operatorname{spec}(J)
=
+0{,}0009
\quad
\text{à }g=0{,}132.
\]

La localisation retenue est donc

\[
\boxed{g_H=0{,}1285}
\qquad
\texttt{[PROP-NUM]}.
\]

La pulsation au franchissement vaut environ

\[
\omega_H\simeq0{,}334,
\]

soit une période proche de

\[
\frac{2\pi}{\omega_H}\simeq18{,}8.
\]

La racine résistante persiste au-delà du seuil. À \(g=0{,}145\), on
trouve encore

\[
x_{\mathrm{res}}=-0{,}2453,
\]

mais l'équilibre est instable. Le calcul valide donc une
\emph{rupture de stabilité}, non une disparition de la branche.

\section{Validation du développement asymptotique}
\label{annexe-num-asymptotique-resistance}

Le développement analytique est

\[
x_{\mathrm{res}}(g)
=
-2{,}1774(g-g_c)
+
0{,}3220(g-g_c)^2
+
O\left((g-g_c)^3\right).
\]

Il est comparé à la racine numérique aux multiples de \(g_c\).

\[
\begin{array}{c|cc}
g & \text{développement} & \text{racine numérique}\\
\hline
2g_c & -0{,}0709 & -0{,}0710\\
3g_c & -0{,}1411 & -0{,}1421
\end{array}
\]

L'accord confirme que la quasi-linéarité observée sur la branche ne
vient pas d'un simple ajustement visuel. Elle résulte de la faible
valeur du terme quadratique après compensation de ses différentes
contributions.

\section{Durée minimale d'appropriation sous R0}
\label{annexe-num-tstar-r0}

Pour une intensité \(g\), la durée critique \(T^\ast(g)\) est repérée
en faisant varier la durée d'une exposition rectangulaire, puis en
classant la trajectoire après découplage.

Le prédicteur linéaire non autonome est comparé à la dyade complète :

\[
\begin{array}{c|cc}
g & T^\ast_{\text{prédit}} & T^\ast_{\text{mesuré}}\\
\hline
0{,}05 & +\infty & +\infty\\
0{,}08 & 31{,}6 & 31{,}7\\
0{,}10 & 24{,}8 & 24{,}8\\
0{,}15 & 17{,}8 & 17{,}9\\
0{,}20 & 14{,}9 & 15{,}0\\
0{,}30 & 12{,}2 & 12{,}3
\end{array}
\]

L'écart relatif reste inférieur à \(0{,}5\,\%\) pour les valeurs
finies testées. La zone de résistance est correctement classée comme
non appropriante, même lorsque la durée tend vers l'infini.

Le prédicteur n'est pas ajusté sur ce tableau. Il utilise la charge
exponentielle de la trace, la stabilité instantanée du repos et la
position de la selle autonome.

\section{Contrôle par le protocole de montée}
\label{annexe-num-rampes}

La dépendance du seuil de capture au protocole est testée à cible
identique :

\[
g=0{,}10.
\]

Trois montées sont comparées : contact brutal, rampe courte de durée
\(20\), et rampe lente de durée \(400\).

Le contact brutal et la rampe courte conduisent à l'appropriation. La
rampe lente laisse la trajectoire suivre la branche résistante stable
et conduit à la résistance.

Ce contrôle montre que

\[
g_{\mathrm{ign}}
=
g_{\mathrm{ign}}(x_0,v_0,\theta_0;\mathcal P)
\]

est une frontière de capture dépendante de l'histoire de chargement.
Il ne s'agit pas d'un seuil structurel universel.

\section{Calcul du cycle R2}
\label{annexe-num-cycle-r2}

Pour chaque valeur de \(g>g_K\), le cycle événementiel de R2 est
calculé en deux étapes.

\subsection{Période de l'interface}

La période \(\Delta(g)\) est la solution positive de l'équation
transcendante

\[
p\left(\Delta;\chi_\ast^+(\Delta),r_P\Theta\right)=\Theta.
\]

Une résolution scalaire fournit \(\Delta\), puis
\(\chi_\ast^-\), \(\chi_\ast^+\), \(\theta_\infty^+\) et
\(\theta_\infty^-\).

\subsection{Vérification de la trace moyenne}

La formule analytique est

\[
\bar\theta
=
\frac{\kappa_Q\chi_\ast^-}{\beta\Delta}.
\]

Quatre contrôles donnent :

\[
\begin{array}{c|cc}
 & \text{prédite} & \text{mesurée}\\
\hline
1&0{,}1293&0{,}1288\\
2&0{,}2886&0{,}2885\\
3&0{,}3953&0{,}3952\\
4&0{,}7216&0{,}7216
\end{array}
\]

L'accord atteint \(10^{-4}\) sur les trois dernières valeurs. Le
premier écart reste compatible avec les erreurs de régime transitoire,
de localisation des événements ou de moyenne sur une durée finie.

\section{Multiplicateur de Floquet}
\label{annexe-num-floquet}

La trace périodique \(\theta_\ast(t)\) produit une matrice linéaire
périodique \(A_g(t)\) pour le bloc \((x,v)\).

La matrice fondamentale est intégrée sur une période :

\[
\dot\Phi_g(t)=A_g(t)\Phi_g(t),
\qquad
\Phi_g(0)=I_2.
\]

La relaxation de la vitesse au Kairos est ajoutée par

\[
R_V=
\begin{pmatrix}
1&0\\
0&r_V
\end{pmatrix}.
\]

La monodromie est

\[
\mathcal M_g=R_V\Phi_g(\Delta(g)).
\]

Le multiplicateur dominant est

\[
\mu_{\max}(g)=\rho_{\mathrm{sp}}(\mathcal M_g).
\]

Le seuil \(g_F\) est repéré par

\[
\mu_{\max}(g_F)=1.
\]

Le calcul donne

\[
\boxed{g_F=0{,}0763}
\qquad
\texttt{[PROP-NUM]}.
\]

Le critère quasi-statique

\[
\bar\theta(g_F^{qs})=\theta_{\mathrm{ign}}
\]

donne

\[
\boxed{g_F^{qs}=0{,}0760.}
\]

L'écart relatif est d'environ \(0{,}4\,\%\). Ce nombre confirme que la
moyenne de la trace capture presque entièrement le mécanisme local,
parce que la période du cycle, environ \(3\) à \(4\), est plus courte
que le temps d'oubli

\[
\frac{1}{\beta}\simeq8{,}3.
\]

Le symbole \(g_F^{qs}\) reste une approximation de \(g_F\), et non un
seuil supplémentaire.

\section{Rampe sur la résistance pulsée}
\label{annexe-num-rampe-resistance-pulsee}

Pour rechercher une rupture de la résistance sous R2, la simulation
est initialisée sur l'attracteur résistant pulsé, puis \(g\) est
augmenté lentement.

Aucune perte de stabilité n'est observée jusqu'à

\[
g=0{,}30.
\]

Le résultat doit être consigné sous la forme

\[
\boxed{g_H^{R2}>0{,}30}
\]

et non comme une localisation de seuil.

Ce contrôle ne permet pas de distinguer un seuil plus élevé, une
résistance stable sur tout le domaine, une autre bifurcation ou un
changement d'attracteur suivi silencieusement par la rampe.

\section{Protocole des deux expositions}
\label{annexe-num-deux-expositions}

Les bandes de sommation sont étudiées avec

\[
g=0{,}18,
\qquad
T_{\mathrm{sub}}=41{,}9,
\qquad
D\in\{2,4,\ldots,60\}.
\]

La première exposition ne produit pas d'appropriation. Elle laisse un
état résiduel qui évolue librement pendant \(D\). Une seconde
exposition est ensuite appliquée.

L'issue est classée après découplage : \(S\) si la trajectoire rejoint
le puits positif, et échec sinon.

Le balayage produit la signature :

\begin{verbatim}
D :  2 ......................................60
     [S.....SSSS....SSS....SSS......]
\end{verbatim}

Les bandes ont une période proche de \(14\).

\section{Spirale libre et période prédite}
\label{annexe-num-spirale-libre}

Après la décroissance suffisante de la trace, le repos autonome est une
spirale amortie. La linéarisation vérifie

\[
\ddot x+\gamma_0\dot x+(1-\zeta)x=0.
\]

La pulsation amortie est

\[
\boxed{
\omega
=
\sqrt{(1-\zeta)-\frac{\gamma_0^2}{4}}
}.
\]

Avec les paramètres de référence,

\[
\boxed{
\frac{2\pi}{\omega}=14{,}63.
}
\]

Cette période est prédite indépendamment du balayage des délais. Elle
explique l'espacement des bandes mesurées.

\section{État de découplage et extrapolation libre}
\label{annexe-num-extrapolation-libre}

L'état mesuré à la fin de la première exposition est

\[
x_d=-0{,}0061,
\qquad
v_d=-0{,}0011,
\qquad
\theta_d=0{,}7255.
\]

Pendant le délai,

\[
\theta(t)=\theta_de^{-\beta t}.
\]

La trajectoire libre linéarisée \(x_{\mathrm{lib}}(t)\) est obtenue en
intégrant

\[
\ddot x_{\mathrm{lib}}
+
\Gamma\!\left(\sqrt2\,\theta_de^{-\beta t}\right)
\dot x_{\mathrm{lib}}
+
\left[
m\!\left(\sqrt2\,\theta_de^{-\beta t}\right)-\zeta
\right]
x_{\mathrm{lib}}
=0,
\]

avec

\[
x_{\mathrm{lib}}(0)=x_d,
\qquad
\dot x_{\mathrm{lib}}(0)=v_d.
\]

Cette extrapolation utilise l'état mesuré une seule fois à la fin de
la première exposition. Elle n'est pas ajustée cellule par cellule.

\section{Temps de recharge avant la seconde instabilité}
\label{annexe-num-temps-charge}

La seconde exposition ne déstabilise pas instantanément le repos.
L'interface doit d'abord reconstruire une charge suffisante.

Le temps de recharge est mesuré indépendamment sur l'application de
retour :

\[
\boxed{t_c=4{,}74.}
\]

La phase pertinente n'est donc pas celle de l'état au début de la
seconde exposition, mais celle qu'il aurait au moment où l'instabilité
devient active : \(D+t_c\).

Le critère testé est

\[
\boxed{
\text{sommation}
\quad\Longleftrightarrow\quad
x_{\mathrm{lib}}(D+t_c)>0.
}
\]

Aucune constante n'est ajustée sur le balayage des délais : \(\omega\)
vient de la linéarisation autonome, l'état initial vient du découplage
mesuré et \(t_c\) est mesuré séparément sur la carte.

\section{Performance du prédicteur de phase}
\label{annexe-num-performance-phase}

Dans le régime

\[
10\leq D\leq48,
\]

le prédicteur obtient

\[
\boxed{18\text{ classifications correctes sur }20.}
\]

Les deux erreurs se trouvent aux bords de bande. Elles correspondent à
un décalage d'environ \(2\) unités de délai, soit la résolution du
balayage.

Tous les cas classés comme sommants rejoignent

\[
x\simeq+2{,}30.
\]

Les échecs de phase relaxent vers

\[
x\simeq0
\]

en passant par la résistance pulsée.

Une fois l'instabilité ouverte, le taux de croissance est proche de
\(0{,}2\). Sur la durée de la seconde exposition, l'amplification est
de l'ordre de

\[
e^8.
\]

Dans ce régime, l'amplitude résiduelle décide donc moins que son signe.
Le côté positif est amplifié vers le puits positif ; le côté négatif
est capturé par la résistance.

\section{Trois régimes de mémoire temporelle}
\label{annexe-num-trois-regimes}

Le balayage distingue trois régimes.

\subsection{Mémoire de charge}

Pour

\[
D\lesssim10,
\]

les variables \(p\), \(\chi\) et \(\theta\) restent fortement
chargées. À \(D=2\), par exemple,

\[
\theta\simeq0{,}72.
\]

Cette charge peut produire une sommation même lorsque la phase
mécanique est défavorable.

\subsection{Mémoire de phase}

Pour

\[
10\lesssim D\lesssim48,
\]

la charge résiduelle a suffisamment décru pour que le signe de

\[
x_{\mathrm{lib}}(D+t_c)
\]

devienne le principal discriminant.

\subsection{Structure fine près de l'oubli}

Pour

\[
D\gtrsim50,
\]

l'amplitude résiduelle entre dans une zone où la structure fine des
bassins redevient déterminante. Le critère de signe seul n'est plus
suffisant.

La forme générique

\[
A(D)\cos(\omega D+\varphi_0)>A_{\mathrm{crit}}
\]

doit donc être lue ainsi : dans le régime de phase,
\(A_{\mathrm{crit}}\simeq0\) ; près de l'extinction, l'amplitude et la
structure fine des bassins ne peuvent plus être négligées.

\section{Contrôle à charge livrée comparable}
\label{annexe-num-controle-charge}

Pour vérifier que les bandes ne sont pas simplement dues à une dose
totale différente, on calcule

\[
Q_{\mathrm{net}}
=
\sum_n
\kappa_Q\chi_n^-.
\]

Deux délais donnent des charges presque identiques :

\[
D=14
\quad\Longrightarrow\quad
Q_{\mathrm{net}}=8{,}255,
\]

\[
D=26
\quad\Longrightarrow\quad
Q_{\mathrm{net}}=8{,}322.
\]

Pourtant, \(D=14\) conduit à l'appropriation, tandis que \(D=26\)
conduit à l'oubli.

\begin{quote}
\textbf{Contrôle.}
Une dose livrée presque identique peut produire deux issues opposées.
La quantité totale de trace injectée ne suffit donc pas à expliquer la
sommation. La phase de réception est nécessaire.
\end{quote}

\section{Amorçage scalaire de l'interface}
\label{annexe-num-amorcage-scalaire}

Une variable lente scalaire est ajoutée :

\[
\dot s=\varepsilon_s(\chi-s),
\]

avec

\[
\varepsilon_s=0{,}005\ll\lambda_\chi.
\]

Le gain des quanta devient

\[
\boxed{
\kappa_Q(s)=\kappa_Q(1+c_ss),
}
\]

avec

\[
c_s=8.
\]

Une première exposition sous-critique de durée \(41{,}9\) est suivie
d'un délai d'oubli de \(300\).

Après ce délai :

\[
|x|\simeq10^{-18},
\qquad
\theta\simeq10^{-15},
\qquad
\chi\simeq10^{-28},
\qquad
p\simeq10^{-17},
\]

mais

\[
\boxed{s=0{,}0095.}
\]

Toutes les variables rapides sont revenues à leur base, tandis que la
variable lente d'interface conserve une trace. Le gain initial des
quanta est alors augmenté de

\[
c_ss=0{,}076,
\]

soit \(7{,}6\,\%\).

\section{Effet de l'amorçage}
\label{annexe-num-effet-amorcage}

La seconde exposition est appliquée à

\[
g=0{,}18.
\]

Les durées \(T_2\) sont balayées dans une fenêtre où la phase produit
des bandes d'appropriation.

L'expérience initiale compare deux états strictement identiques, sauf
pour la valeur de \(s\) au début de la seconde exposition :

\[
\begin{array}{c|cc}
 & \text{cellules appropriées} & \text{première durée appropriante}\\
\hline
s_0=0 & 17/24 & T_2\simeq75\\
s_0=0{,}0095 & 24/24 & T_2\leq40
\end{array}
\]

La première exposition n'a changé aucun bassin autonome, mais elle a
modifié l'accessibilité future de ces bassins :

\[
\boxed{T_2^\ast<T_1^\ast.}
\]

Le seuil soutenu varie très peu :

\[
\Delta g_{\mathrm{app}}=-0{,}0014.
\]

La mémoire d'interface change donc surtout ce qui est accessible dans
la durée finie d'une rencontre, et non ce qui devient possible après
une exposition infiniment longue.

\section{Contrôle : gain constant contre boucle dynamique}
\label{annexe-num-controle-gain-boucle}

L'effet précédent peut venir d'une simple tête d'avance due au gain
initial de \(7{,}6\,\%\), ou d'une boucle dynamique dans laquelle des
quanta plus grands accélèrent les variables qui augmenteront les
quanta suivants.

Trois bras sont comparés pour \(T_2=40,\ldots,155\).

\[
\begin{array}{l|cc}
\text{bras} & \text{cellules appropriées} & \text{première durée}\\
\hline
\text{sans amorçage} & 17 & 75\\
\text{gain constant }(+7{,}6\,\%) & 23 & 45\\
\text{boucle dynamique complète} & 24 & 40
\end{array}
\]

Le résultat impose une interprétation prudente :

\[
\boxed{
\text{l'essentiel de l'amorçage vient de la tête d'avance.}
}
\]

La boucle dynamique ajoute une cellule appropriante et réduit la
première durée de \(5\) unités. Elle existe, mais elle est secondaire
pour le paramétrage testé.

\section{Contrôle du signe scalaire}
\label{annexe-num-controle-signe-scalaire}

Une valeur négative de \(c_s\) devait tester une éventuelle
habituation.

Avec

\[
c_s=-8
\]

sans garde de positivité, certains quanta deviennent négatifs. Le
système entre alors dans un régime d'anti-écriture qui efface
directement la trace. Ce comportement est un artefact de la loi
scalaire, pas une habituation interprétable.

Avec

\[
c_s=-2,
\]

le levier résiduel est d'environ \(-1{,}9\,\%\). Aucun effet robuste
n'est détecté.

Le contrôle négatif ne valide donc pas une habituation scalaire. Il
montre seulement qu'une comparaison équitable demanderait un levier
négatif de même ordre que le levier positif, avec une loi maintenant
les quanta dans un domaine admissible.

\section{Mémoire vectorielle et compatibilité}
\label{annexe-num-memoire-vectorielle}

La mémoire scalaire \(s\) retient une intensité d'exposition, mais pas
sa direction.

Pour tester une mémoire spécifique au contenu reçu, on utilise une
variable vectorielle

\[
\boldsymbol s\in\mathbb R^2.
\]

La règle d'écriture au Kairos devient

\[
\boxed{
\tau^+
=
\tau^-
+
\max\left(
1+c_s\left\langle
\boldsymbol s,
\widehat{\boldsymbol\chi}
\right\rangle,
0
\right)
\kappa_Q\boldsymbol\chi^-.
}
\]

Le facteur

\[
\left\langle
\boldsymbol s,
\widehat{\boldsymbol\chi}
\right\rangle
\]

mesure la compatibilité entre la mémoire laissée par les expositions
passées et le contenu actuellement présenté.

\section{Protocole de spécificité}
\label{annexe-num-protocole-specificite}

Le modèle est simulé en dimension deux avec le préconditionneur complet
\(M(\tau)\).

Le protocole est :

\begin{enumerate}
\item première exposition dans la direction \(u_1\) ;
\item délai d'oubli de \(300\) ;
\item seconde exposition dans l'une des directions
  \[
  u_1,
  \qquad
  u_\perp,
  \qquad
  -u_1;
  \]
\item balayage \(T_2=40,\ldots,150\).
\end{enumerate}

La petite perturbation initiale est appariée au début de la seconde
exposition dans tous les bras.

Cette précaution est nécessaire : sans elle, l'érosion de la
perturbation pendant le long délai d'oubli peut masquer la mémoire
d'interface et produire une différence qui vient seulement de
l'amplitude initiale disponible.

\section{Résultats de spécificité}
\label{annexe-num-resultats-specificite}

Le tableau suivant indique le nombre de durées appropriantes parmi
douze valeurs testées. Le nombre entre parenthèses est la première
durée appropriante.

\[
\begin{array}{l|cccc}
\text{mémoire}
&
\text{même direction}
&
\text{orthogonale}
&
\text{opposée}
&
\text{sans amorce}
\\
\hline
\text{scalaire}
&5/12\;(110)
&5/12\;(110)
&5/12\;(110)
&2/12\;(140)
\\
\text{vectorielle}
&5/12\;(110)
&2/12\;(140)
&1/12\;(140)
&2/12\;(140)
\end{array}
\]

Trois conclusions sont permises.

\subsection{La mémoire scalaire produit une sensibilisation globale}

La facilitation est identique pour les trois directions. La variable
scalaire mémorise qu'une exposition a eu lieu, mais pas son contenu.

\subsection{La mémoire vectorielle est directionnelle}

Pour la même direction, la facilitation est conservée :

\[
5/12>2/12.
\]

Pour la direction orthogonale,

\[
\left\langle\boldsymbol s,u_\perp\right\rangle=0.
\]

L'amorce devient exactement invisible dans la règle de compatibilité :

\[
2/12=2/12
\]

par rapport au bras sans amorce.

\subsection{L'opposition produit une inhibition}

Pour la direction opposée,

\[
\left\langle\boldsymbol s,-u_1\right\rangle<0.
\]

Le gain des quanta diminue et le nombre de durées appropriantes devient

\[
1/12<2/12.
\]

La même loi géométrique produit donc

\[
\boxed{
\text{facilitation}
\quad\text{ou}\quad
\text{inhibition},
}
\]

selon le signe de la compatibilité.

\section{Interprétation permise}
\label{annexe-num-interpretation-specificite}

Le résultat démontre numériquement une \emph{mémoire d'interface
spécifique au contenu reçu} ou une \emph{mémoire directionnellement
individualisée}.

Il ne démontre pas encore une mémoire de l'identité du partenaire. Deux partenaires distincts produisant la même direction projetée restent indiscernables pour le protocole considéré.

Il reste à tester plusieurs partenaires de directions voisines, la transférabilité d'un amorçage, l'oubli différentiel des composantes de \(\boldsymbol s\), la saturation de la mémoire vectorielle et l'effet de l'ordre de plusieurs amorçages.

\section{Synthèse des contrôles}
\label{annexe-num-synthese-controles}

\begin{center}
\begin{tabular}{p{0.25\linewidth}p{0.31\linewidth}p{0.32\linewidth}}
\hline
\text{Question}
&
\text{Contrôle}
&
\text{Conclusion}
\\
\hline
Les bandes viennent-elles de la dose ?
&
\(D=14\) et \(D=26\), charges \(8{,}255\) et \(8{,}322\)
&
Non : issues opposées malgré une dose comparable.
\\
L'amorçage vient-il d'une boucle cumulative ?
&
Gain constant contre gain dynamique
&
Principalement de la tête d'avance ; boucle secondaire.
\\
La mémoire scalaire code-t-elle un contenu ?
&
Même, orthogonale et opposée
&
Non : facilitation identique.
\\
La mémoire vectorielle est-elle sélective ?
&
Produit scalaire nul ou négatif
&
Oui : neutralité orthogonale et inhibition opposée.
\\
Une valeur négative de \(c_s\) suffit-elle à produire une habituation
scalaire ?
&
\(c_s=-8\) puis \(c_s=-2\)
&
Non concluante : anti-écriture ou effet trop faible.
\\
\hline
\end{tabular}
\end{center}

\section{Hiérarchie des résultats numériques}
\label{annexe-num-hierarchie-resultats}

Les résultats du Tome II s'organisent en quatre niveaux.

\subsection{Structure locale}

\[
g_c=0{,}0327
\]

pour le modèle continu, et

\[
g_F=0{,}0763
\]

pour le cycle R2 réduit.

\subsection{Activité de l'interface}

\[
g_K=0{,}03608
\]

marque l'apparition d'une activité événementielle soutenue.

\subsection{Franchissement de bassin}

\[
g_{\mathrm{app}}^{R2}(x_0;\mathcal P)=0{,}1632
\]

est une valeur protocolaire.

\subsection{Résistance pulsée}

\[
g_H^{R2}>0{,}30
\]

est une borne expérimentale, non un seuil localisé.

La hiérarchie correcte est

\[
\boxed{
g_K<g_F<g_{\mathrm{app}}^{R2}<g_H^{R2}.
}
\]

\section{Sources d'incertitude}
\label{annexe-num-incertitudes}

Les principales sources d'incertitude sont :

\begin{itemize}
\item la résolution du balayage en \(g\), \(T\) ou \(D\) ;
\item la localisation numérique des Kairos ;
\item la durée finie utilisée pour classer une trajectoire ;
\item la proximité d'une séparatrice ;
\item l'érosion d'une petite perturbation pendant un long délai ;
\item la sensibilité des bords de bande à la non-linéarité de la capture.
\end{itemize}

Une valeur de seuil obtenue sur une grille est donc accompagnée de son
statut : formule, racine numérique, borne ou valeur protocolaire. Cette
distinction est plus importante qu'un nombre excessif de décimales.

\section{Tests de robustesse recommandés}
\label{annexe-num-tests-recommandes}

Avant de figer une version destinée à diffusion, les scripts doivent
permettre les tests suivants :

\begin{enumerate}
\item division par deux des tolérances du solveur ;
\item raffinement local du maillage autour des seuils ;
\item doublement de la durée autonome utilisée pour classifier l'issue ;
\item déplacement contrôlé de la perturbation initiale ;
\item répétition avec une localisation d'événements plus stricte ;
\item vérification des identités analytiques
  \[
  \bar\theta=\frac{\kappa_Q\chi_\ast^-}{\beta\Delta}
  \]
  et
  \[
  p(\Delta)=\Theta;
  \]
\item conservation des mêmes perturbations dans les comparaisons appariées ;
\item export du rapport complet en format structuré.
\end{enumerate}

Ces tests ne changent pas le mécanisme étudié. Ils vérifient que la
conclusion ne dépend pas d'un détail de discrétisation.

\section{Questions ouvertes}
\label{annexe-num-questions-ouvertes}

\subsection{Bords des bandes}

Le prédicteur de signe commet deux erreurs sur vingt, toutes deux aux
bords. Il reste à relier quantitativement le décalage de phase à la
non-linéarité de la capture.

\subsection{Habituation scalaire}

Une loi scalaire doit encore être construite pour produire un levier
négatif comparable au levier positif sans autoriser des quanta
pathologiques.

\subsection{Rupture de la résistance pulsée}

La valeur exacte de \(g_H^{R2}\), si elle existe, n'est pas connue.

\subsection{Mémoire du partenaire}

La mémoire vectorielle encode une direction projetée, pas une identité.
Il faut introduire plusieurs partenaires ou plusieurs canaux pour
tester cette distinction.

\subsection{Validation conditionnelle complète}

La règle de compatibilité vectorielle modifie le quantum lors du saut.
Elle n'a pas encore été intégrée à toute la théorie de R2. Les valeurs

\[
g_F,
\qquad
g_{\mathrm{app}}^{R2},
\qquad
g_H^{R2}
\]

devront être recalculées sous cette loi.

\subsection*{Raccord avec le Tome III}

La question de l'ordre des histoires n'est plus laissée ouverte ici.
Le Tome III distingue trois mécanismes susceptibles de produire une
différence entre deux ordres d'épisodes : la pondération de récence,
la conditionnalité de l'écriture et le chevauchement des transitoires.
La présente annexe reste consacrée au soutien numérique des mécanismes
des Tomes I et II ; l'analyse de la composition des histoires relève
désormais du Tome III.

\section{Conclusion}
\label{annexe-num-conclusion}

Les expériences numériques établissent une chaîne cohérente :

\[
\boxed{
\begin{minipage}{0.86\linewidth}
Le paysage autonome fournit les attracteurs et les séparatrices.
L'interface transforme l'amplitude, la cadence et la phase des
influences qui donnent accès à ces bassins. Deux expositions de dose
comparable peuvent produire des issues différentes parce qu'elles
rencontrent le système à des phases différentes. Une première
exposition non appropriée peut ensuite modifier durablement
l'accessibilité future des bassins. Dans la version vectorielle, cette
modification devient sélective : la même compatibilité géométrique
produit facilitation, neutralité ou inhibition.
\end{minipage}
}
\]

Ce résultat ne constitue pas une preuve globale de l'individuation. Il
fournit les briques numériques nécessaires à son étude :

\[
\text{appropriation des bassins}
\quad+\quad
\text{amorçage directionnel des accessibilités}.
\]
